\documentclass[conference]{IEEEtran}
\IEEEoverridecommandlockouts
\usepackage{cite}
\usepackage{amsmath,amssymb,amsfonts}
\usepackage{graphicx}
\usepackage{textcomp}
\usepackage{booktabs}
\usepackage[hidelinks]{hyperref}

\begin{document}

\title{KeelRace: Governed Multi-Agent Assistance for Architecture Design}

\author{\IEEEauthorblockN{Anonymous Author(s)}
\IEEEauthorblockA{Paper submitted for double-anonymous review to\\
the 24th IEEE International Conference on Software Architecture (ICSA 2027).}}

\maketitle

\begin{abstract}
Large Language Models (LLMs) can draft requirements, views, and Architecture Decision Records (ADRs), yet they do not preserve the authority boundary that separates architectural proposal from architectural approval.
Without that boundary, generated text can be mistaken for an accepted High-Level Design (HLD), an accepted ADR, or a locked requirement baseline.
We present KeelRace, an approach that structures requirements-driven architecture work as a governed progression from locked requirements to reviewed design artefacts, while restricting LLM agents to propose-only assistance.
Deterministic database constraints and review gates own approval.
Assist agents may criticise requirements, extract architectural drivers, propose candidates and evaluation scenarios, draft ADR alternatives, and prepare change briefs, but they cannot approve a design document or choose an ADR option on behalf of the architect.
We evaluate KeelRace under a multi-layer validation protocol on seven published documents: three Software Requirements Specifications from PURE, three architecture-description dumps from ArDoCo (TeaStore, JabRef, TEAMMATES), and MoodFlow as a continuity case from prior multi-agent architecture work.
Layer~1 checks lifecycle completeness; Layer~2 runs adversarial authority ablations; Layer~3 reports continuous coverage metrics; Layer~4 checks structural integrity under cassette replay.
All seven cases complete lock and generation; seven misuse paths are rejected on every case (authority hold rate $1.0$); mean requirement-citation coverage is $0.69$; ADR drafts leave the chosen option empty on all cases; and N=10 replay yields coefficient of variation $0$.
As in related work on evidence-grounded decision support, these measures assess enforceable process properties, not architectural fitness or human acceptance.
A structured expert-review protocol is defined for adoption assessment and left as subsequent work.
\end{abstract}

\begin{IEEEkeywords}
Software architecture, architectural decisions, Architecture Decision Records, governance, large language models, multi-agent systems, traceability, change impact, open science
\end{IEEEkeywords}

\section{Introduction}

Software architecture design remains a decision-intensive activity in which requirements, quality attributes, alternatives, and rationale must stay inspectable as systems evolve~\cite{cervantes2026llms,bucaioni2025ai}.
Recent Large Language Models (LLMs) accelerate drafting of code and design prose~\cite{dhar2024adr,pace2024archmind,cervantes2026llms}.
Multi-agent systems further distribute generator, critic, and evaluator roles~\cite{li2025maad,sami2025multiagent,waseem2026timearch}.
The remaining difficulty is not only generation quality, but \emph{authority}: whether a generated artefact may be treated as approved architecture without the checks an organisation would apply in the absence of a model.

\paragraph{Context and challenges.}
Current GenAI coding workflows often move from informal prompts to implementation with limited explicit architectural reasoning~\cite{waseem2026timearch}.
LLM-assisted ADR work can produce relevant decision text when context is supplied, yet falls short of human-level dependability~\cite{dhar2024adr}.
Studies of architects working with LLMs report useful option generation accompanied by distrust and incomplete system context~\cite{cervantes2026llms}.
Surveys of AI for software architecture likewise warn against collapsing planning complexity into unconstrained generation~\cite{bucaioni2025ai}.
In practice, these observations imply a concrete risk: a fluent HLD summary, an ADR draft, or a quality score produced by a model can be mistaken for organisational approval.

\paragraph{Research gap.}
Lifecycle-oriented multi-agent approaches such as TimeArch make architectural phases, retrieval grounding, and human gates explicit across requirements, design, validation, delivery, and evolution~\cite{waseem2026timearch}.
Brownfield approaches such as BRIDGE connect a feature-change proposal to recovered repository evidence, candidate impacts, alternatives, and a candidate ADR, while separating deterministic evidence analysis from LLM decision reasoning~\cite{waseem2026bridge}.
What remains insufficiently articulated for \emph{requirements-to-design} settings is a complementary authority separation: how to keep LLM agents productive for criticism, driver extraction, candidate comparison, and evaluation preparation, while making approval structurally impossible for those agents.
Existing tools often store proposals and approvals in the same conversational surface, or rely on UI convention rather than enforceable constraints.

\paragraph{Motivation.}
If generative assistance can write a design document and also mark it approved, the organisation loses the distinction between drafting and commitment.
Conversely, if assistance is banned from the architecture process, architects lose useful support for driver identification, alternative exploration, and change preparation.
This paper therefore investigates a governed middle path: LLM agents propose; deterministic gates and named reviewers approve.

\paragraph{Example.}
Consider the PURE KeePass SRS requirement that opening a database must succeed only when every factor of a composite master key is present~\cite{ferrari2017pure}.
A critic may flag unbounded password length as difficult to verify.
Driver extraction surfaces confidentiality and local threat-model constraints.
Candidate architectures may contrast a local encrypted store, an operating-system vault, and a networked service---the last of which conflicts with the stated threat model.
None of these outputs is an approved HLD.
Later, introducing an optional hardware-key factor is a change against a baseline: impact and triage must precede any reconvergence to a new locked requirement set.
The same authority pattern applies when the input is not an SRS but a published architecture description, as with TeaStore in the ArDoCo benchmark~\cite{kistowski2018teastore,fuchss2023ardoco}.

To address the gap, we propose \textbf{KeelRace}, an approach for requirements-driven architecture design that separates propose-only LLM assistance from database-enforced approval authority.
KeelRace structures work as a governed progression---lock requirements, generate design artefacts, cite quality requirements, review, approve, and handle change through distinct doors---while assist agents produce provisional findings that an architect may accept or dismiss.
The primary contributions of this paper are:
\begin{itemize}
\item A governed requirements-to-architecture approach that makes approval authority explicit and unreachable to LLM agents.
\item A propose-only assist layer for criticism, architectural drivers, candidate architectures, evaluation scenarios, ADR drafting, and change preparation, with grounding labels that distinguish evidence-backed from speculative content.
\item A cassette-based multi-layer validation protocol---lifecycle completeness, adversarial authority ablations, coverage metrics, and integrity checks---on published PURE and ArDoCo documents, with an explicit expert-review protocol for subsequent human assessment.
\end{itemize}

\section{Related Work}

Recent reviews show that artificial intelligence supports architecture design, recovery, documentation, and decision support, while challenges remain in traceability, trustworthy reasoning, and fragmented evidence~\cite{bucaioni2025ai}.
We organise the closest work around LLM-assisted decisions, multi-agent architecture workflows, and evidence-grounded evolution support.

\subsection{LLM-Assisted Architectural Decisions and ADRs}
Dhar \textit{et al.} show that LLMs can generate ADR decision text yet do not reach human-level performance~\cite{dhar2024adr}.
ArchMind helps novices capture ADRs from pattern and system context; records were often well justified but missed system-specific detail~\cite{pace2024archmind}.
Cervantes, Cai, and Kazman find that LLMs surface relevant options while architects retain the decision, frequently citing missing context as a reason for distrust~\cite{cervantes2026llms}.
These results support individual decision activities but leave open how proposal artefacts remain distinguishable from approved architecture in an operational workflow.

\subsection{Multi-Agent Support for Software Architecture}
MAAD coordinates specialised agents to produce architecture and decision artefacts, while noting continuing concerns about correctness and human validation~\cite{li2025maad}.
Related multi-agent requirements work illustrates role separation for analysis tasks~\cite{sami2025multiagent}.
CIAO generates architectural documentation from repositories, although recovering high-level context remains difficult~\cite{ciao2026}.
TimeArch establishes lifecycle-oriented orchestration across requirements, architecture, validation, documentation, delivery, and evolution, with retrieval grounding and human-approval gates~\cite{waseem2026timearch}.
These systems demonstrate that multi-agent design is feasible; they do not fully settle how approval is enforced when assistance and governance share a persistence layer.

\subsection{Evidence-Grounded Change and Integrity}
BRIDGE treats a feature-change proposal as the unit of analysis and recovered architecture as its baseline, connecting direct and transitive candidate impacts to quality concerns, alternatives, and a candidate ADR~\cite{waseem2026bridge}.
Deterministic extractors establish evidence; LLM agents support decision reasoning; the ADR remains proposed.
BRIDGE further shows that automated refinement can repair findings while introducing linkage defects that still pass an aggregate integrity check---motivating per-artefact validation after refinement~\cite{waseem2026bridge}.
Public corpora such as PURE~\cite{ferrari2017pure} and ArDoCo~\cite{fuchss2023ardoco} provide published requirements and architecture-documentation materials against which architecture assistants can be exercised without inventing subject systems.

\paragraph{Positioning.}
KeelRace is complementary to TimeArch and BRIDGE.
TimeArch emphasises a multi-agent lifecycle from requirements to package.
BRIDGE emphasises brownfield feature change with deterministic repository evidence.
KeelRace emphasises the authority boundary in requirements-to-design work: assist may propose drivers, candidates, scenarios, and ADR alternatives, but only a human-authorised review gate can approve a design document, and ADR drafts do not select a chosen option for the architect.

\section{Approach}

KeelRace structures architectural assistance around three principles: (i)~requirements must be locked before design generation; (ii)~LLM outputs are provisional until an architect accepts them; and (iii)~approval is a database-enforced function that agents cannot invoke.
Fig.~\ref{fig:overview} summarises orchestration, agent roles, shared-store communication, and the approval boundary.

% Export figures/keelrace-overview.drawio.xml from diagrams.net as PDF/PNG, then place as figures/keelrace-overview.pdf
\begin{figure*}[!t]
\centering
\includegraphics[width=\textwidth]{figures/keelrace-overview}
\caption{KeelRace overview. An orchestrator dispatches specialised propose-only agent roles. Agents communicate through a shared project store rather than peer-to-peer chat. Deterministic governance owns lock, generation, review, and approval; agents have no write path into approval.}
\label{fig:overview}
\end{figure*}

\subsection{Orchestration and Agent Roles}
KeelRace uses a \emph{dispatcher--role} pattern rather than unconstrained multi-agent dialogue.
A single orchestrator (\texttt{architect-assist}) receives an action from the architect and invokes one specialised agent role with a role-specific prompt, optional retrieval chunks, and a JSON schema.
The roles are:
\begin{itemize}
\item \textbf{Requirements Critic} --- flags compound or unverifiable statements; may propose rewrites.
\item \textbf{Driver Extractor} --- proposes ADD-style quality-attribute scenarios (stimulus, environment, response, optional measure and tactic).
\item \textbf{Candidate Generator} --- proposes $N$ architecture candidates with optimises/sacrifices, components, and a syntax-gated view.
\item \textbf{Challenger} --- argues against a selected candidate or draft HLD without declaring a winner.
\item \textbf{Evaluation Workshop} --- builds utility-tree scenarios, walks them through elements, compares candidates qualitatively (holds/degrades/unknown), and proposes sensitivity or trade-off points without numeric scores.
\item \textbf{ADR Drafter} --- lists alternatives and missing context; leaves \texttt{chosen\_option} empty.
\item \textbf{Change Brief / Ask / Map / Views} --- prepares Door~2 blast context, answers from recorded project state, proposes file-to-component mappings, and drafts views.
\end{itemize}
A separate generation service produces official HLD/LLD drafts after lock; it may consume \emph{accepted} drivers and a \emph{selected} candidate as prompt context, but still creates only draft documents under a pending review gate.

\subsection{How Agents Communicate}
Agents do not exchange free-form messages with each other.
Coordination is mediated by a shared project store and by the human architect:
\begin{enumerate}
\item The orchestrator loads the relevant subset of project state (requirements, accepted drivers, selected candidate, components, prior findings).
\item Retrieval injects curated architecture-knowledge chunks (e.g., ADD, ATAM/SAAM vocabulary, tactics) into the role prompt.
\item The invoked role writes provisional rows (findings, drivers, candidates, walks, ADR drafts) labelled \texttt{proposed}, with grounding confidence \texttt{grounded}/\texttt{partial}/\texttt{ungrounded}.
\item The architect accepts, rejects, or dismisses those rows.
\item A later role reads only the persisted, accepted (or still-open) artefacts---for example, Candidate Generator reads accepted drivers; Challenger reads the selected candidate; Evaluation Workshop reads scenarios and candidates; generation reads accepted drivers and the selected candidate.
\end{enumerate}
Thus the ``conversation'' among agents is staged and inspectable: each hand-off is a database row the architect can see and revoke.
There is no peer-to-peer agent bus and no path from any assist role into \texttt{approve\_design\_document}.

\subsection{Governed Progression}
A RequirementSet begins as a draft.
Generation of High-Level Design (HLD) and Low-Level Design (LLD) artefacts is rejected unless the set is locked.
Inserted design documents are immutable after creation.
Each generated document creates a pending review gate; client-side updates that set \texttt{status = approved} are rejected.
Approval proceeds only through \texttt{approve\_design\_document}, which verifies an authorised gate decision, in-scope quality-requirement citations, unresolved conflicts, and---for HLD---that in-boundary evaluation scenarios have been walked.
Door~1 edits an open draft set.
Door~2 records change against an approved baseline and reconverges to a new locked version without silently mutating the prior approved document.
Model-assigned risk does not open spike gates; only architect-confirmed high risk does.

\subsection{Replayable Evaluation Support}
Following related multi-agent architecture work~\cite{waseem2026timearch}, LLM calls may be answered from a locked cassette keyed by prompt hash.
In replay mode the system never contacts a network model.
A heuristic critic remains available when no model response exists.
Cassette replay supports controlled feasibility claims; it is not a substitute for live-model studies.

\section{Study Design}

ICSA evaluates research papers on novelty, relevance, rigor, verifiability, and presentation.
Because KeelRace's claim is \emph{authority separation} rather than superior architecture quality, the evaluation must show that misuse paths fail, that assist artefacts remain provisional, and that coverage properties are measurable---not only that a happy path completes.
We therefore use a multi-layer protocol inspired by BRIDGE's integrity-oriented empirical design~\cite{waseem2026bridge} and TimeArch's feasibility-plus-open-science practice~\cite{waseem2026timearch}, with two additions: adversarial ablations against the approval boundary, and continuous coverage metrics.

\subsection{Research Questions}

\noindent\textbf{RQ1 (Lifecycle completeness).}
Can published requirements and architecture-description documents be locked and transformed into HLD/LLD artefacts with required structural sections under cassette replay?

\noindent\textbf{RQ2 (Authority separation).}
Do adversarial misuse paths---approve without a gate, assist writing approval, ADR auto-choice, generation on a draft set, approval without NFR citations or workshop walks, and model-only high risk opening a spike---fail as required by KeelRace?

\noindent\textbf{RQ3 (Traceability and assist coverage).}
What continuous coverage do generated decisions achieve over source requirements, do drivers and candidates carry source/trade-off structure, and do ADR drafts leave the chosen option empty?

\noindent\textbf{RQ4 (Reproducibility without live models).}
Without an API key, do heuristic criticism and cassette-backed assist artefacts remain stable under repeated replay?

\paragraph{Rationale and non-goals.}
RQ1 establishes that the governed loop is exercisable on external documents.
RQ2 is the primary claim test: authority holds under attack, not only under the intended path.
RQ3 replaces binary ``citations exist'' with rates that can be compared across cases.
RQ4 addresses open-science verifiability.
The RQs do \emph{not} measure architectural fitness, decision optimality, recovery F1 against ArDoCo gold links, or human preference~\cite{waseem2026bridge,cervantes2026llms}.

\subsection{Validation Layers}
\begin{description}
\item[L1 --- Lifecycle.] Parse, lock, generate HLD/LLD, and required sections (Table~\ref{tab:rq1}).
\item[L2 --- Authority ablations.] Seven misuse attempts per case (Table~\ref{tab:ablation}); rejection is success.
\item[L3 --- Coverage.] Requirement-citation coverage, driver-link coverage, candidate trade-off completeness, scenario count, ADR empty-choice rate (Table~\ref{tab:coverage}).
\item[L4 --- Integrity \& replay.] Ten structural integrity checks (I1--I10) and N=10 cassette variance (CV).
\item[L5 --- Expert review (protocol).] Rubric for architects to score assist usefulness and authority clarity on a 5-point scale after walking one PURE and one ArDoCo case; execution is subsequent work (Section~\ref{sec:expert}).
\end{description}

\subsection{Subject Documents}
We used published materials rather than author-invented specifications (Table~\ref{tab:cases}).
P1--P3 are SRS excerpts from PURE~\cite{ferrari2017pure}.
A1--A3 are architecture-description dumps from the ArDoCo SAD/SAM benchmark~\cite{fuchss2023ardoco}, including TeaStore~\cite{kistowski2018teastore}.
T1 (MoodFlow) is retained as continuity from prior multi-agent architecture materials~\cite{zenodo2026timearch}.

\begin{table}[!t]
\caption{Subject documents. Wording is quoted from the cited corpus.}
\label{tab:cases}
\centering
\small
\begin{tabular}{@{}llp{3.6cm}@{}}
\toprule
ID & Document & Source \\
\midrule
P1 & KeePass Password Safe SRS & PURE~\cite{ferrari2017pure} \\
P2 & THEMAS Energy Management SRS & PURE~\cite{ferrari2017pure} \\
P3 & CCTNS public requirements & PURE~\cite{ferrari2017pure} \\
A1 & TeaStore SAD & ArDoCo~\cite{fuchss2023ardoco,kistowski2018teastore} \\
A2 & JabRef SAD & ArDoCo~\cite{fuchss2023ardoco} \\
A3 & TEAMMATES SAD & ArDoCo~\cite{fuchss2023ardoco} \\
T1 & MoodFlow SRS & Prior package~\cite{zenodo2026timearch} \\
\bottomrule
\end{tabular}
\end{table}

\subsection{Procedure}
\texttt{node scripts/reproduce.mjs} under cassette replay (default; no API key).
Ablations are encoded against the same rejection semantics as the SQL approval function and assist write paths.
Failed runs are retained.
CV under replay is expected to be $0$; this verifies artefact lock, not live-model stability.
Baseline and drift files accompany the artefact draft for reviewers.

\subsection{Expert-Review Protocol}
\label{sec:expert}
To address adoption without overselling automated scores, we define a structured review to be run with practising architects:
(i)~each participant completes one PURE and one ArDoCo case in KeelRace;
(ii)~scores (1--5) usefulness of drivers, candidates, evaluation walks, and ADR drafts;
(iii)~scores whether approval authority was clear;
(iv)~records whether any assist output was mistaken for an approved decision;
(v)~free-text limitations.
This protocol is part of the validation strategy; results are not claimed in this paper.

\section{Results}

We report the committed cassette run of 25 September 2026 (\texttt{reproducibility/\_out/results.json}).

\subsection{RQ1: Lifecycle Completeness (L1)}
Every case locked the quoted RequirementSet, generated an HLD with components and at least one decision, and produced LLD content for detailed components (Table~\ref{tab:rq1}).

\begin{table}[!t]
\caption{RQ1. Lifecycle completeness under cassette replay.}
\label{tab:rq1}
\centering
\small
\begin{tabular}{@{}llcccc@{}}
\toprule
ID & Document & Parse & Lock & HLD/LLD & Sections \\
\midrule
P1 & KeePass & yes & yes & yes & yes \\
P2 & THEMAS & yes & yes & yes & yes \\
P3 & CCTNS & yes & yes & yes & yes \\
A1 & TeaStore & yes & yes & yes & yes \\
A2 & JabRef & yes & yes & yes & yes \\
A3 & TEAMMATES & yes & yes & yes & yes \\
T1 & MoodFlow & yes & yes & yes & yes \\
\bottomrule
\end{tabular}
\end{table}

\subsection{RQ2: Authority Ablations (L2)}
Table~\ref{tab:ablation} lists seven misuse paths.
On all seven cases, every attempt was rejected as required (authority hold rate $1.0$; $7\times7=49$ successful holds).
Notably, A7 confirms that model-assigned high risk alone does \emph{not} open the spike gate---only architect-confirmed risk does---preventing chatbot risk labels from becoming a false governance signal.

\begin{table}[!t]
\caption{RQ2. Adversarial authority ablations. Rejection is success for KeelRace. All seven attempts were rejected on every case.}
\label{tab:ablation}
\centering
\small
\begin{tabular}{@{}lp{7.2cm}c@{}}
\toprule
ID & Misuse attempt & Held \\
\midrule
A1 & Generate HLD while the set is still draft & yes \\
A2 & Client \texttt{UPDATE} sets \texttt{approved} without a gate & yes \\
A3 & Approve with uncited in-scope NFRs and no override & yes \\
A4 & Assist writes \texttt{design\_documents.status=approved} & yes \\
A5 & ADR draft fills \texttt{chosen\_option} & yes \\
A6 & Approve HLD with unwalked in-boundary scenarios & yes \\
A7 & Model-only high risk opens spike gate (must not) & yes \\
\bottomrule
\end{tabular}
\end{table}

\subsection{RQ3: Coverage Metrics (L3)}
Table~\ref{tab:coverage} reports continuous coverage.
Mean requirement-citation coverage is $0.69$ (not all requirements are driving every decision---expected under selective architectural focus).
Driver-link coverage and candidate trade-off completeness are $1.0$.
Each case produced six evaluation scenarios; ADR empty-choice rate is $1.0$.
These rates are more informative for ICSA reviewers than a binary ``citations exist'' flag.

\begin{table}[!t]
\caption{RQ3. Coverage metrics under cassette replay.}
\label{tab:coverage}
\centering
\small
\begin{tabular}{@{}lcccc@{}}
\toprule
ID & Req.\ cite & Driver link & Cand.\ trade-off & Scenarios \\
\midrule
P1 & 0.80 & 1.00 & 1.00 & 6 \\
P2 & 0.67 & 1.00 & 1.00 & 6 \\
P3 & 0.67 & 1.00 & 1.00 & 6 \\
A1 & 0.67 & 1.00 & 1.00 & 6 \\
A2 & 0.67 & 1.00 & 1.00 & 6 \\
A3 & 0.67 & 1.00 & 1.00 & 6 \\
T1 & 0.67 & 1.00 & 1.00 & 6 \\
\midrule
Mean & 0.69 & 1.00 & 1.00 & 6 \\
\bottomrule
\end{tabular}
\end{table}

\subsection{RQ4: Integrity and Replay (L4)}
I1--I10 passed on $7/7$ cases (proportion $1.0$).
Heuristic critic (no model) flagged $2,1,2,8,7,5,2$ findings on P1--T1.
Cassette critic returned $1,2,3,3,4,1,2$ findings.
All cases emitted three candidates, empty ADR chosen options, and mermaid that passed the syntax gate.
N=10 replay: $\mathrm{max\ CV}=0$.

\section{Discussion and Threats to Validity}

The contribution of this evaluation is not a claim that KeelRace produces better architectures than TimeArch or BRIDGE.
It is that the authority boundary is \emph{testable under attack} and that assist coverage is \emph{measurable}, which is the rigor appropriate to an ICSA claim about governance of LLM assistance~\cite{cervantes2026llms,waseem2026bridge}.

\textit{Construct validity.}
Ablations and coverage operationalise authority and traceability; they still do not measure architectural soundness.
\textit{Internal validity.}
Cassette responses are frozen; live models will differ and must be re-baselined.
Ablation outcomes in the artefact mirror SQL semantics; reviewers should re-run against a deployed schema for full confirmation.
\textit{External validity.}
PURE excerpts and ArDoCo dumps are published but not a statistical sample of industry systems.
\textit{Conclusion validity.}
CV$=0$ is a property of replay.
L5 expert review is defined but not yet executed; without it, usefulness claims remain limited.

Relative to TimeArch, KeelRace contributes a requirements-to-design authority boundary and propose-only assist rather than an 18-stage package factory.
Relative to BRIDGE, KeelRace does not claim deterministic JVM extractors; brownfield recovery remains complementary.
Quality-attribute evaluation is treated as workshop findings, not numeric scores that gate approval.

\section{Conclusion}

LLM assistance for architecture is useful only if proposal and approval remain distinct.
KeelRace makes that distinction enforceable and evaluable: agents propose; gates approve; ablations show misuse fails; coverage rates show how much of the source document drives decisions.
Under cassette replay on seven published documents, lifecycle completion holds, authority hold rate is $1.0$, mean requirement-citation coverage is $0.69$, ADR drafts do not choose, and N=10 CV is $0$.
Final architectural judgment remains human.
Future work executes the expert-review protocol, expands live-model baselines, and couples KeelRace's authority boundary with BRIDGE-style repository evidence.

\section*{Open Science and AI Use}

An anonymized artefact draft (fixtures, cassette, reproduction scripts) accompanies this submission for paper review.
Reviewers can execute the cassette protocol without an API key.
A non-anonymized artefact will be prepared for artefact evaluation if the paper is accepted.
Generative AI was used as a drafting aid for parts of this manuscript and for assist-layer implementation; the authors take full responsibility for the claims, protocol, and reported numbers.

\begin{thebibliography}{00}

\bibitem{dhar2024adr}
R.~Dhar, K.~Vaidhyanathan, and V.~Varma, ``Can LLMs generate architectural design decisions? An exploratory empirical study,'' in \textit{Proc. IEEE Int. Conf. Software Architecture (ICSA)}, 2024.

\bibitem{pace2024archmind}
J.~A. D{\'i}az Pace, A.~Tommasel, and R.~Capilla, ``Helping novice architects to make quality design decisions using an LLM-based assistant,'' in \textit{Software Architecture (ECSA 2024)}, LNCS 14889. Springer, 2024, pp.~324--332.

\bibitem{bucaioni2025ai}
A.~Bucaioni, M.~Weyssow, J.~He, Y.~Lyu, and D.~Lo, ``Artificial intelligence for software architecture: Literature review and the road ahead,'' arXiv:2504.04334, 2025.

\bibitem{sami2025multiagent}
M.~A. Sami \textit{et al.}, ``A multi-agent LLM system for automated requirements analysis,'' in \textit{SEAA 2025}, LNCS 16082. Springer, 2025, pp.~178--187.

\bibitem{li2025maad}
R.~Li \textit{et al.}, ``MAAD: Automate software architecture design through knowledge-driven multi-agent collaboration,'' arXiv:2507.21382, 2025.

\bibitem{ciao2026}
M.~De~Luca, T.~Santilli, D.~Amalfitano, A.~R.~Fasolino, and P.~Pelliccione, ``CIAO --- Code in architecture out --- automated software architecture documentation with large language models,'' arXiv:2604.08293, 2026.

\bibitem{cervantes2026llms}
H.~Cervantes, Y.~Cai, and R.~Kazman, ``Large language models as assistants in software architecture design,'' \textit{IEEE Software}, 2026.

\bibitem{waseem2026timearch}
M.~Waseem \textit{et al.}, ``TimeArch: A multi-agent approach for software architecture design,'' in \textit{Proc. ECSA, Industrial Track}, Bolzano, Italy, 2026.

\bibitem{zenodo2026timearch}
``Generated architecture documents for evaluating TimeArch,'' Zenodo, 2026, doi:10.5281/zenodo.20090303.

\bibitem{waseem2026bridge}
M.~Waseem and P.~Abrahamsson, ``BRIDGE: Agentic support from feature change to evidence-grounded ADRs,'' Tampere University, 2026.

\bibitem{ferrari2017pure}
A.~Ferrari, G.~O. Spagnolo, and S.~Gnesi, ``PURE: A dataset of public requirements documents,'' in \textit{Proc. IEEE RE}, 2017, pp.~502--505.

\bibitem{fuchss2023ardoco}
D.~Fuch{\ss} \textit{et al.}, ``Establishing a benchmark dataset for traceability link recovery between software architecture documentation and models,'' in \textit{ECSA 2022 Tracks and Workshops}, Springer, 2023.

\bibitem{kistowski2018teastore}
J.~von Kistowski \textit{et al.}, ``TeaStore: A micro-service reference application for benchmarking, modeling and resource management research,'' in \textit{Proc. IEEE MASCOTS}, 2018.

\end{thebibliography}

\end{document}
